\scanpage{125}

% manuscript lines / position: top centre
\mspage{465}

% manuscript lines / position: 1–10
Il me semble que notre époque se caractérise comme celle d'une exacerbati\unsure{on} à outrance de cette dégradation culture\deleted{\unsure{bt|b}}\unsure{ll}e. Parmi les derniers actes de cette histoire, il y a ceux, intimement solidaires, de la ``course à l'espace'' entre le\unsure{s} deux superpuissances antagonistes (imbues de valeurs essentiellement identiques) et de la course aux armements (nucléaires notamment). Comme acte ultime et dénou\unsure{e}ment probable de cette évolution forcenée dans les surenchères d'un certain type de ``force'' ou de ``pouvoir'', on peut prévoir dès à présent quelque holocauste nucléaire (ou autre, il y a l'embarras du choix...) à l'échelle planétaire. Il aur\unsure{a} peut-être ce mérite de résoudre tous les problèmes d'un seul coup et une fois po\unsure{ur} toutes...

% manuscript lines / position: 11–25
Mon propos ici n'est pas pourtant de brosser un alléchant tableau de ``fin du monde'' (on ne m'a pas attendu pour cela), et encore moins de partir en guerre contre le ``muscle'', ou contre ``le cerveau'' (alias l'``esprit\unsure{''}). Je sais bien que même mes ``tripes'' n'auraient rien à y gagner! Je tiens à mes muscles et à mon cerveau, qui me sont bien utiles on s'en doute, comme je tiens aussi à mes ``tripes'', qui ne le sont pas moins. Plutôt, il me semble utile de dire ici en quelqu\unsure{es} mots (si faire se peut) comment s'est joué dans ma propre personne ce conflit pr\unsure{o}fond, véhiculé par la culture environnante, entre ces deux types de valeurs. En termes plus terre à terre, il s'agit aussi de l'historique de mes attitudes (d'acceptation voire d'exaltation, ou de rejet) de deux \mbox{a s p e c t s} ou \mbox{f a c e s} également réels et tangibles de ma personne, inséparables et complémentaires par nature\handwritten{\unsure{,}} et nullement antagonistes par eux-mêmes. Je pourrais les appeler ``\mbox{l' h o m m e}'' et ``\mbox{l a\quad f e m m e}'' en moi, ou aussi (pour prendre des appe\handwritten{\unsure{ll}}ations moins ``chargées'', et qui pour cela offrent moins de risque d'induire en erreur), le ``\mbox{y a n g}'' et le ``\mbox{y i n}''.

% manuscript lines / position: 26–35
Il semblerait que chez la plupart des personnes, les ``jeux sont faits'' dès la petite enfance, où se mettent en place les mécanismes essentiels qui, pendant la vie entière, vont dominer en silence, avec une efficacité d'automate parfaitement au point, attitudes et comportements. Au coeur de\unsure{.} ces mécanismes sont ceux d'affirmation ou de rejet de tels et tels traits en nous, ou de telles pulsions profondes, à ``signature'' soit yang soit yin, ou de tels \unsure{et tels} ``paquets'' de traits et de pulsions à signature donnée, voire même du paquet ``yang'' ou du paquet ``yin'' tout entier. Ce sont ces mécanismes qui, dans une très large mesure, déterminent tous les autres mécanismes de choix (affirmation ou rejet) structura\unsure{nt} notre ``moi''.

% manuscript lines / position: bottom left
% archive sheet number: 124
